\chapter{Estado da arte}
\label{cap:contexto}

Este capítulo apresenta o enquadramento científico e tecnológico do trabalho. Caracteriza-se
primeiro o público a que o sistema se dirige e as limitações que definem o problema; examinam-se
depois as soluções atualmente disponíveis; percorrem-se por fim as áreas técnicas de onde provêm
os componentes utilizados, situando a contribuição desta dissertação em relação a cada uma. Quem o leia fica a saber que as três perguntas do investidor já têm resposta na literatura, e por que razão nenhuma ferramenta gratuita as junta.

\section{O investidor particular e o problema da atenção}
\label{sec:ctx_investidor}

Esta secção caracteriza o comportamento documentado do público a que o sistema se dirige, uma vez
que são as limitações desse comportamento, e não a dimensão do mercado, que delimitam o problema
tratado.

O Capítulo~\ref{cap:introducao} estabeleceu a escala do universo em causa. A caracterização que
falta é de comportamento, e começa por uma observação que contraria a representação corrente do
investidor particular como agente que apenas reage em pânico. Durante a queda de março de 2020, os
investidores particulares de uma plataforma sem comissões aumentaram as suas posições agregadas em
vez de reduzirem a exposição \autocite{welch2022robinhood}. Trata-se de um público que decide com
regularidade, e que por isso beneficia de melhor informação no momento da decisão.

A literatura de finanças comportamentais estabelece que essas decisões se afastam de forma
sistemática do ideal racional \autocite{barberis2003behavioral}. Duas regularidades condicionam
diretamente o desenho de um sistema de alertas. A primeira é a aversão à perda
\autocite{kahneman1979prospect}, já enunciada no capítulo anterior, e cuja consequência para este
trabalho é que o custo da ausência de contexto é assimétrico: é maior precisamente quando a
informação é desfavorável.

A segunda é o efeito da atenção, e é a que fundamenta o problema. \textcite{barber2008glitters}
estabelecem que os investidores particulares são compradores líquidos das ações que lhes captam a
atenção, identificadas por três indicadores observáveis, que são a presença em notícias, o volume
de transação anormal e o retorno extremo num único dia. O mecanismo proposto pelos autores é o do
problema de procura: perante um número muito elevado de alternativas, o investidor não avalia todas
e a sua escolha recai sobre o subconjunto que lhe foi tornado visível.

Os autores delimitam o
alcance da afirmação, que respeita a um desequilíbrio agregado entre compras e vendas, sem
abranger um comportamento universal, e reportam a conclusão que justifica este trabalho: o padrão de aquisição
assim formado não gera retornos superiores. A base empírica são registos de corretoras
norte-americanas entre 1991 e 1999, anteriores à negociação por telemóvel e à ausência de
comissões, pelo que o que se assume transferível é o mecanismo e não os valores medidos. A
centralidade da atenção neste processo é de tal ordem que se tornou mensurável através do volume de
pesquisas na internet \autocite{da2011attention}.

Os três indicadores identificados pela literatura são calculados por este sistema, mas o estatuto
de cada um difere. A notícia e o movimento extremo constituem os dois gatilhos descritos no
Capítulo~\ref{cap:introducao} e podem, por si, originar um alerta. O volume é calculado e
apresentado como contexto, e não interrompe o utilizador. A razão desta assimetria é empírica e
está reportada na Secção~\ref{sec:av_sintese}: a combinação do volume com os restantes sinais
melhorou a triagem em apenas um de três orçamentos considerados, resultado que não sustenta a
introdução de uma terceira via de interrupção. A ideia de fundir sinais de origens distintas num
indicador único provém de um painel de informação em tempo real com essa arquitetura
\autocite{worldmonitor2026}, sugerido pelo coorientador deste trabalho, e entra na dissertação na
condição em que foi avaliada, ou seja como alternativa medida e não adotada.

Há, sobre este ponto, evidência de campo mais recente e mais próxima do objeto deste trabalho, e
que aponta em sentido desfavorável. \textcite{liu2023alerts} acompanharam 20\,904 investidores
particulares de uma plataforma de fundos de investimento em Taiwan, entre janeiro de 2017 e junho
de 2020, dos quais 932 ativaram os alertas de preço da própria plataforma. Estimando o efeito por
diferenças em diferenças, quantificam-no numa degradação do desempenho da carteira da ordem de um
ponto percentual em seis meses. Os dois mecanismos que identificam são o aumento de cerca de 7\% no
número de transações e o facto de os fundos que o investidor mexeu depois da adesão terem
desempenho pior do que os que não mexeu. A leitura dos autores é a de que o alerta dirige a atenção
para a volatilidade de curto prazo e induz negociação em excesso.

O resultado é incómodo para uma dissertação que constrói um sistema de alertas, e é por isso que
consta do corpo do texto e não de uma nota. Convém, ainda assim, ser exato quanto ao que foi
medido. O alerta estudado dispara quando o preço de um fundo atravessa um limiar definido pelo
próprio investidor e comunica-lhe o preço atual, sem qualquer indicação sobre a causa; é o alerta
de limiar que a Secção~\ref{sec:ctx_ferramentas} volta a encontrar nas aplicações das corretoras. O
mecanismo de dano que os autores descrevem, que é a negociação por realimentação do preço em vez
do uso de informação sobre o valor da empresa, é precisamente o que se instala quando ao
destinatário se entrega um número e mais nada. Os próprios autores concluem que melhorar o acesso à
informação não basta, e que é preciso atender também à capacidade de a processar.

Esta dissertação toma o resultado como razão de desenho e não como validação do desenho que já
tinha. Dele decorrem duas opções. A primeira é que o alerta transporta a evidência que o justifica,
em vez de comunicar apenas que um limiar foi atravessado. A segunda é que a quantidade de alertas é
limitada por decisão explícita, e não pelo que os dados produzirem, conforme a
Secção~\ref{sec:sis_politica}. Se estas duas opções bastam para evitar o efeito que
\textcite{liu2023alerts} mediram é matéria por demonstrar, e o Capítulo~\ref{cap:conclusoes}
retoma-a nesses termos.

Do lado das instituições financeiras, a adoção de inteligência artificial deixou de ser
experimental, como documenta o inquérito a empresas do setor conduzido pelo
\textcite{ccaf2026aifs}. Esse inquérito não mede a população dos produtos destinados ao investidor
particular, pelo que não sustenta afirmações sobre ela; a caracterização dessa população assenta no
que os próprios fornecedores declaram, e é objeto da secção seguinte.

\section{Ferramentas de informação financeira para o retalho}
\label{sec:ctx_ferramentas}

Esta secção examina as soluções atualmente disponíveis, avaliando cada uma contra um critério fixo,
que são as três perguntas enunciadas no Capítulo~\ref{cap:introducao}. A
Tabela~\ref{tab:ctx_ferramentas} apresenta o resultado dessa avaliação.

\begin{table}[!htbp]
\caption[As ferramentas existentes contra as três perguntas]{As categorias de ferramenta existentes
avaliadas contra as três perguntas do Capítulo~\ref{cap:introducao}. A classificação
\emph{parcial} indica que a ferramenta fornece matéria-prima para a resposta e deixa a
interpretação ao utilizador.}
\label{tab:ctx_ferramentas}
\centering
\small
\begin{tabular}{L{0.22\textwidth} C{0.11\textwidth} C{0.15\textwidth} C{0.12\textwidth}
                L{0.17\textwidth}}
\toprule
\tabhead{Ferramenta} & \tabhead{É invulgar?} & \tabhead{Empresa ou mercado?} &
\tabhead{Já aconteceu?} & \tabhead{Custo} \\
\midrule
Aplicação da corretora & não & não & não & incluído \\
Agregador de notícias & não & não & não & gratuito \\
Aplicação de sentimento & parcial & não & não & gratuito ou baixo \\
Gestor automático de carteira & não & não & não & percentagem da carteira \\
Terminal profissional & sim & sim & sim & não publicado \\
\midrule
\textbf{Este trabalho} & \textbf{sim} & \textbf{sim} & \textbf{sim} & \textbf{gratuito} \\
\bottomrule
\end{tabular}
\end{table}

A tabela evidencia o essencial. A única ferramenta existente que responde às três perguntas é
também a mais dispendiosa de todas. O problema não é, portanto, de natureza científica, mas de acesso: as respostas são
conhecidas e estão condicionadas a uma subscrição incompatível com o perfil de investidor em causa.

Três categorias merecem exame individual, uma vez que cada uma falha por uma razão distinta e essas
razões informam decisões tomadas nos capítulos seguintes.

A aplicação da corretora é o instrumento de acesso mais imediato ao preço, e é aquele cuja
documentação pública descreve o menor conjunto de indicadores. Apresenta a variação percentual do
dia e um gráfico de preço. Não indica se essa variação é elevada
para aquela ação em particular, não separa a componente de empresa da componente de mercado e não
dispõe de memória de acontecimentos comparáveis.

Os alertas de preço destas aplicações partilham a limitação de forma mais acentuada. Disparam
quando um limiar fixo é atravessado, critério que ignora a volatilidade própria de cada ação e
produz por isso alertas frequentes nas empresas agitadas e raros nas calmas. Comunicam ainda que um nível foi
ultrapassado, e nunca porquê. Esta observação determina a opção metodológica descrita na
Secção~\ref{sec:met_detecao}, que substitui o limiar absoluto por uma medida relativa à norma
recente de cada empresa, e cuja diferença é quantificada na Secção~\ref{sec:av_qi1}.

As aplicações de sentimento classificam notícias como positivas ou negativas e apresentam o
resultado dessa classificação. O que fornecem é uma leitura da notícia e não do movimento: uma
classificação fortemente negativa indica que algo terá acontecido, e não estabelece se a
variação do preço é elevada para aquela ação, que é a primeira pergunta. É neste sentido
restrito que a Tabela~\ref{tab:ctx_ferramentas} lhes atribui uma resposta parcial, e não
avançam além dela. A Secção~\ref{sec:av_qi2} apresenta ainda uma medição própria segundo a qual o tema de uma
notícia e a direção do movimento que se lhe segue coincidem com menos frequência do que a intuição
sugere, o que limita o alcance de uma pontuação de sentimento como indicador de impacto.

Os gestores automáticos de carteira respondem sobretudo a uma pergunta distinta, a da distribuição
de capital por um conjunto de ativos. A literatura reconhece-lhes benefícios genuínos, embora não
uniformes: \textcite{dacunto2019robo} medem uma melhoria da diversificação e uma redução da
volatilidade nos investidores que detinham menos de cinco ações antes da adesão, ao passo que para
os que já detinham mais de dez o efeito sobre a diversificação é próximo de nulo. A contenção de
enviesamentos comportamentais é o benefício mais consistentemente reportado
\autocite{dacunto2019robo, cardillo2024robo}. A diferença face a este trabalho é de objeto e não de
transparência: um gestor automático recomenda uma carteira e executa-a, ao passo que este sistema
não formula recomendações e entrega evidência sobre um acontecimento concreto. Acresce que a
prestação de aconselhamento sobre investimentos constitui atividade regulada, fronteira que o
presente trabalho evita por desenho, conforme discutido na Secção~\ref{sec:met_etica}.

Existem, por fim, produtos recentes que declaram responder exatamente à pergunta tratada nesta
dissertação. O Robinhood Cortex, anunciado em março de 2025, declara na página do próprio
fornecedor destinar-se a responder à questão de saber por que razão uma ação está a subir ou a
descer num dado dia \autocite{robinhood2025cortex}. Os momentos-chave do Google Finance,
apresentados em junho de 2026, propõem-se explicar por que razão uma ação se moveu
\autocite{google2026finance}. Trata-se da mesma pergunta, formulada por organizações com recursos e
base de utilizadores incomparavelmente maiores.

A diferença situa-se no que as respetivas páginas
declaram ser entregue. Ambas descrevem um resumo em linguagem natural, e nenhuma descreve a
apresentação dos casos passados individualizados, com datas, semelhanças medidas e comportamento
observado do preço após cada um, que é o que este trabalho entrega. A distinção reivindicada não é
de fluência, mas de verificabilidade, e respeita ao que as páginas declaram: os produtos não foram
testados, pelo que nada se afirma sobre o que efetivamente entregam.

Até aqui esta secção examinou produtos. Falta a literatura académica, e omiti-la deixaria a
afirmação de lacuna assente na inspeção de páginas de fornecedores, que é base frágil. A classe de
sistema aqui em causa, que é vigiar acontecimentos e notificar o investidor particular, é tratada
em investigação publicada há cerca de duas décadas, e o exemplo mais desenvolvido é o \emph{MoFiN
DSS} de \textcite{muntermann2009ubiquitous}. O paralelo é próximo ao ponto de ser desconfortável:
também se dirige ao investidor particular, também parte da observação de que este reage às
notícias com mais intensidade e menos rapidez do que o investidor institucional, também é
construído e avaliado segundo o paradigma de investigação por desenho que a
Secção~\ref{sec:met_metodologia} adota, e também publica a avaliação num artigo de revista.

Vale a pena precisar o que esse sistema faz, porque é dessa precisão que sai o posicionamento
deste trabalho. Sobre 425 comunicados de empresas alemãs publicados entre 2003 e 2005, e séries de
preços intradiárias ao minuto, o autor ajusta por regressão dois modelos: um estima o efeito de
preço que se seguirá ao comunicado, e o outro estima durante quanto tempo esse efeito se
mantém. O comunicado é notificado ao investidor quando o efeito estimado excede o efeito médio
observado. A avaliação é uma simulação \emph{ex ante} do excedente do consumidor, ou seja mede a
oportunidade económica que a notificação abre, e não a qualidade daquilo que é comunicado. O
utilizador recebe a indicação de que um comunicado é relevante e não recebe a razão pela qual o é.

Três consequências para esta dissertação. A primeira é que a lacuna não pode ser enunciada como
ausência de sistemas de alerta para o investidor particular, porque eles existem e estão
publicados; face a este sistema, distingue-se pela postura epistémica, entre estimar o que vai
acontecer e explicar o que já aconteceu, e como diferença naquilo que é entregue com o alerta. A
segunda é que a medida de relevância de \textcite{muntermann2009ubiquitous} é um retorno anormal
\emph{em valor absoluto}, ajustado pelo mercado e corrigido pelo comportamento habitual da própria
ação, ou seja magnitude e não direção. É, no essencial, a mesma grandeza que a
Secção~\ref{sec:met_detecao} mede, com duas diferenças que interessam: aqui é medida depois do
facto, e serve para explicar em vez de antecipar. Que um trabalho publicado em 2009 tenha chegado
à magnitude, e não à direção, como definição operacional de acontecimento relevante para este
público é um apoio independente a uma opção que esta dissertação tomou por outra via. A terceira
é que aquele trabalho avaliou o artefacto por simulação e remeteu explicitamente para investigação
futura o estudo com utilizadores reais, exatamente a limitação que o
Capítulo~\ref{cap:conclusoes} enuncia a respeito deste. Não é atenuante, mas situa o problema:
trata-se de uma dificuldade da área e não de uma omissão isolada.

Outros trabalhos aproximam-se da explicação do movimento observado. O \emph{Squawk Bot} de
\textcite{dang2020squawk} aprende conjuntamente séries de preços e texto para selecionar notícias
associadas ao comportamento da cotação. O \emph{DeepTrust} de \textcite{chan2022deeptrust}
combina deteção de anomalias, recuperação de publicações do Twitter e avaliação da sua
fiabilidade. Este último é uma dissertação disponibilizada no arXiv, cuja avaliação se concentra
em dois episódios; a sua existência impede que explicar anomalias por recuperação de informação
seja apresentado como novidade deste trabalho.

O \emph{Stock Pattern Assistant} de \textcite{neela2025spa}, também disponibilizado no arXiv,
segmenta preços em sequências monótonas, associa acontecimentos por proximidade temporal e propõe
uma narrativa gerada a partir desses objetos, sem previsão nem atribuição causal. A demonstração
usa acontecimentos de uma amostra construída, pelo que não equivale a uma avaliação contínua
sobre notícias capturadas. Explicar sem prever é, portanto, uma opção partilhada. O posicionamento
desta dissertação assenta na combinação das três perguntas, na evidência entregue e na avaliação
dos componentes em funcionamento, sem reivindicar a prioridade daquela opção.

A objeção mais razoável ao trabalho consiste em perguntar por que razão um assistente de
inteligência artificial genérico não resolveria o mesmo problema. Não foi realizada uma comparação
com nenhum desses assistentes, pelo que o que se segue é um argumento de desenho e não um resultado.
Um assistente que não disponha de acesso externo a dados de mercado não dispõe de três elementos.
O primeiro é a volatilidade recente da ação, que permite determinar se uma variação é elevada. O
segundo é a separação entre a componente de mercado e a componente específica da empresa. O
terceiro é uma base de casos passados sobre a qual efetuar uma recuperação.

Quanto a este último ponto, um modelo de linguagem recorda
casos vistos durante o treino, o que não constitui recuperação e não é verificável pelo leitor.
Existem modelos de grande dimensão especializados em texto financeiro, cuja capacidade está
documentada \autocite{li2023llmfinance}. A saída de qualquer deles permanece uma afirmação que se lê com facilidade e não se confere.

A
literatura de explicabilidade dá razões para desconfiar dela. \textcite{lipton2018mythos}
argumenta que a interpretabilidade não designa uma propriedade única e que um sistema deve
declarar qual delas oferece. \textcite{rudin2019stop} argumenta que uma explicação produzida sobre
um modelo opaco é uma aproximação que pode induzir em erro. Nenhum dos dois trata de texto gerado por modelos de linguagem, pelo que a transposição é do
autor e o que dela se retira é uma razão de desenho e não um resultado da literatura. É esta a razão pela qual a linguagem nunca produz
factos neste sistema: os números provêm dos componentes que os calcularam, e o texto é composto a
partir deles, conforme descrito na Secção~\ref{sec:sis_explicacao}.

\section{Deteção de anomalias em séries temporais financeiras}
\label{sec:ctx_anomalias}

Esta secção situa a primeira decisão do sistema no campo da deteção de anomalias e enuncia as
famílias de método consideradas.

A deteção de anomalias constitui uma área estabelecida, com uma taxonomia consolidada que separa os
métodos pela forma como definem normalidade \autocite{chandola2009anomaly, pang2021deep}. Três
famílias são relevantes para séries temporais financeiras.

A família estatística define normalidade a partir da distribuição local dos próprios dados,
estimada numa janela deslizante, e assinala como anómalas as observações que dela se afastam além
de um limiar. A explicação de uma deteção coincide com o cálculo que a produziu, o que a torna
integralmente comunicável ao utilizador. Em contrapartida, supõe que a distribuição dos retornos é
localmente estável.

A família aprendida não supervisionada define normalidade a partir da estrutura geométrica dos
dados. O Isolation Forest \autocite{liu2008isolation} identifica como anómalos os pontos que uma
partição aleatória isola com poucos cortes; o Local Outlier Factor \autocite{breunig2000lof}
identifica-os pela densidade reduzida face à dos seus vizinhos. Ambos dispensam rótulos e captam
padrões multivariados que uma regra sobre uma única série não capta, ao custo de a decisão deixar
de ser enunciável numa frase.

A família profunda aprende um modelo de normalidade através de uma rede \autocite{pang2021deep}.
Oferece maior capacidade de representação, exige volumes de dados superiores e afasta-se do
requisito de explicabilidade que estrutura este trabalho.

Uma restrição prática condiciona toda esta escolha. Em domínios financeiros vizinhos, como a
deteção de fraude, os métodos não supervisionados ocupam posição central precisamente porque
anomalias rotuladas são escassas e os padrões se alteram ao longo do tempo
\autocite{ahmed2016financial}. A mesma condição verifica-se aqui: não existe um conjunto de dados
em que tenha sido registado, dia a dia, que movimentos justificariam um alerta. Esta ausência
determina a forma como o Capítulo~\ref{cap:avaliacao} pode avaliar o componente, e é a razão pela
qual a avaliação assenta num critério derivado dos próprios dados em vez de um julgamento humano.

A suposição de estabilidade local em que a família estatística assenta é confrontada e não
escondida. Os modelos de volatilidade condicional, introduzidos por \textcite{engle1982arch} e
generalizados por \textcite{bollerslev1986garch}, formalizam o agrupamento de volatilidade, ou seja
a tendência de os períodos calmos e os períodos agitados ocorrerem em séries. Ambos os artigos
formalizam o fenómeno sobre séries de inflação, pelo que a sua transposição para retornos de ações
não é assumida: a Secção~\ref{sec:av_qi1} mede-a nos dados deste trabalho, comparando a janela de
pesos iguais com uma variante que atribui maior peso às observações recentes. Uma janela deslizante
acompanha parcialmente o fenómeno, uma vez que a norma contra a qual cada dia é comparado é
reestimada diariamente a partir do passado próximo. A escolha não elimina o problema, e a medição
referida quantifica o que custa mantê-la.

A comparação entre a família estatística e a família aprendida não permanece no plano do argumento.
O Capítulo~\ref{cap:avaliacao} executa ambas sobre os mesmos dados e reporta o resultado obtido.

\section{Representação de texto e recuperação semântica}
\label{sec:ctx_texto}

Esta secção percorre as formas de representar texto que tornam possível comparar notícias por
significado, e enuncia o critério de avaliação da recuperação que delas decorre.

A comparação de notícias por significado depende inteiramente da representação adotada, e essa
representação atravessou três gerações, esquematizadas na Figura~\ref{fig:ctx_geracoes} e
detalhadas na Tabela~\ref{tab:ctx_texto}.

\begin{figure}[ht]
\centering
\begin{tikzpicture}[
  font=\small, >={Stealth[length=2.4mm]},
  g/.style={draw, rounded corners, align=center, text width=3.0cm, minimum height=3.3cm,
            inner sep=5pt, fill=black!3},
]
  \node[g] (a) {\textbf{Contagem de palavras}\\[3pt]
    \footnotesize cada palavra é\\ \footnotesize um símbolo isolado\\[3pt]
    \footnotesize \emph{``rose'' e ``gained''}\\ \footnotesize \emph{não têm relação}};
  \node[g, right=0.85cm of a] (b) {\textbf{Vetores de palavra}\\[3pt]
    \footnotesize palavras de uso\\ \footnotesize próximo ficam perto\\[3pt]
    \footnotesize \emph{um sentido}\\ \footnotesize \emph{por palavra}};
  \node[g, right=0.85cm of b] (c) {\textbf{Vetores de frase}\\[3pt]
    \footnotesize o sentido depende\\ \footnotesize do contexto\\[3pt]
    \footnotesize \emph{representação}\\ \footnotesize \emph{adotada neste trabalho}};
  \draw[->, thick, black!45] (a) -- (b);
  \draw[->, thick, black!45] (b) -- (c);
\end{tikzpicture}
\caption[Três gerações de representação de texto]{As três gerações de representação de texto. Cada
uma resolve uma limitação identificada na anterior, e a terceira é a que permite comparar duas
frases diretamente pela posição dos respetivos vetores.}
\label{fig:ctx_geracoes}
\end{figure}

A primeira geração representa um documento pela frequência dos termos que contém. O modelo de
espaço vetorial \autocite{salton1975vsm} é transparente e computacionalmente barato, e é cego a
sinónimos: os termos ``rose'' e ``gained'' constituem, para esta representação, símbolos sem
relação. Dentro da mesma geração existe uma segunda linha, oriunda de um quadro probabilístico
distinto, a das funções de ordenação de que o BM25 \autocite{robertson2009bm25} é o representante
corrente; ambas permanecem linhas de base fortes em recuperação de informação e partilham aquela
cegueira. Em texto financeiro existe ainda uma variante especializada, os léxicos de domínio, que
corrigem o facto de os dicionários de sentimento genéricos interpretarem incorretamente palavras
correntes na linguagem financeira \autocite{loughran2011liability}. São integralmente
transparentes, e ficam limitados ao vocabulário que foi compilado.

A segunda geração representa cada palavra como um vetor num espaço onde a proximidade traduz uso
semelhante, aprendido a partir do contexto local \autocite{mikolov2013word2vec} ou da coocorrência
global \autocite{pennington2014glove}. Esta representação aproxima sinónimos, resolvendo a limitação
anterior, e introduz uma nova: cada palavra recebe um vetor único, independentemente da frase em
que ocorre.

A terceira geração remove essa limitação. O Transformer \autocite{vaswani2017attention} tornou
viável o processamento da frase completa com dependências de longo alcance, e o \gls{BERT}
\autocite{devlin2019bert} produziu representações contextuais, nas quais o vetor de uma palavra
depende das que a rodeiam. Um passo adicional é necessário para o problema aqui tratado: comparar
duas frases com um modelo dessa natureza exigiria processá-las conjuntamente, uma vez por par, o
que não é praticável sobre uma base com dezenas de milhares de casos. O \gls{SBERT}
\autocite{reimers2019sbert} resolve esta questão produzindo um vetor por frase, calculado uma única
vez, de modo que a comparação passa a ser uma operação aritmética entre vetores já existentes.

A contribuição do \gls{SBERT} que importa a este trabalho não reside na arquitetura mas no objetivo
de treino: os vetores são aprendidos de forma que a distância entre dois deles traduza semelhança
de sentido. Os autores reportam que um codificador da família \gls{BERT} utilizado diretamente para
comparar frases produz representações fracas para esse fim, ficando por vezes atrás da média
simples de vetores de palavra \autocite{reimers2019sbert}. Daqui decorre que um codificador ajustado
a outra tarefa não herda esta propriedade pelo facto de ter sido treinado no domínio adequado.

\begin{table}[!htbp]
\caption[Abordagens de representação de texto]{As abordagens consideradas para representar notícias
financeiras, com a representação que cada uma produz e as implicações da sua adoção.}
\label{tab:ctx_texto}
\centering
\small
\begin{tabular}{L{0.22\textwidth} L{0.28\textwidth} L{0.36\textwidth}}
\toprule
\tabhead{Abordagem} & \tabhead{Representação} & \tabhead{Forças e limitações} \\
\midrule
Contagem de termos \newline \autocite{salton1975vsm, robertson2009bm25}
  & Frequência ponderada de palavras
  & Transparente e rápida; cega a sinónimos \\
\addlinespace
Léxicos financeiros \newline \autocite{loughran2011liability}
  & Listas de palavras com polaridade atribuída
  & Totalmente transparente; limitada ao vocabulário compilado \\
\addlinespace
Vetores de palavra \newline \autocite{mikolov2013word2vec, pennington2014glove}
  & Um vetor fixo por palavra
  & Aproxima sinónimos; um só sentido por palavra \\
\addlinespace
Vetores de frase \newline \autocite{devlin2019bert, reimers2019sbert}
  & Um vetor por frase, sensível ao contexto
  & Compara frases diretamente; modelo maior e menos legível internamente \\
\addlinespace
Modelos de domínio \newline \autocite{liu2020finbert, huang2023finbert}
  & Como acima, ajustada a texto financeiro
  & Vocabulário do domínio; superioridade nesta tarefa não é automática \\
\addlinespace
Modelos de grande dimensão \newline \autocite{li2023llmfinance}
  & Geração de texto sobre conhecimento amplo
  & Capacidade elevada; dispendiosos, opacos e fora da restrição de gratuitidade \\
\bottomrule
\end{tabular}
\end{table}

Para texto financeiro existem modelos treinados especificamente para o domínio
\autocite{liu2020finbert, huang2023finbert}. A suposição natural
é a de que sejam superiores nesta tarefa por essa razão. A suposição não é automaticamente
verdadeira, e a Secção~\ref{sec:av_qi2} mede-a em vez de a assumir.

O trabalho inscreve-se na literatura mais vasta de processamento de linguagem natural aplicado a
finanças. A revisão de \textcite{kearney2014textual} cobre o estado dessa área até 2014, dominado
pelos léxicos de domínio e por classificadores supervisionados sobre contagens de palavras; as
representações contextuais de frase utilizadas neste trabalho são posteriores e não constam dessa
revisão. Existe ainda uma linha de investigação substancial que utiliza texto para antecipar
movimentos de preço \autocite{xing2018nlffsurvey, ding2015deep}. O presente trabalho situa-se fora
dessa linha por opção de desenho: o texto é aqui empregue para localizar casos anteriores
comparáveis, cujo desfecho já ocorreu e já foi medido, e não para gerar uma previsão.

Um trabalho recente dessa linha merece exame individual, por ser o mais próximo daquilo que aqui
se faz. \textcite{du2024contrastive} treinam uma rede de tripletos sobre mensagens de redes
sociais e séries de preços, e apresentam ao utilizador, junto de cada previsão, o caso histórico
mais semelhante com o mesmo desfecho e o mais semelhante com desfecho oposto. É recuperação de
análogos usada como explicação, sobre os mesmos materiais que este trabalho usa, e é por isso o
vizinho publicado mais próximo.

As diferenças são, ainda assim, de fundo, e vale a pena enunciá-las porque delimitam a
contribuição. O rótulo que organiza os tripletos é a direção do movimento do dia seguinte, e a
tarefa avaliada é prevê-la; a magnitude entra apenas para definir uma faixa morta em torno de
zero. O codificador de frases é usado congelado, como camada de entrada de uma rede convolucional
e de uma rede recorrente com atenção, que são essas sim treinadas, ou seja o espaço em que as
notícias são comparadas não é alterado. E a explicação é demonstrada por três estudos de caso com
mapas de atenção, sem medição da sua qualidade e sem participantes.

Há naquele trabalho um resultado próprio que merece registo, porque diz respeito à dificuldade da
tarefa e não ao método. No estudo de ablação, aplicar a aprendizagem contrastiva apenas à
componente textual degrada a exatidão nos três conjuntos de referência, ficando abaixo da variante
que não usa aprendizagem contrastiva nenhuma; os ganhos vêm da componente quantitativa, e os
próprios autores concluem que a informação textual desempenha um papel complementar. Quem
pretenda usar texto para antecipar a direção de um movimento tem aqui, publicada, uma indicação de
que o sinal é fraco; a Secção~\ref{sec:av_qi2} chega, por via inteiramente independente, a uma
conclusão do mesmo sinal.

É sobre essas duas diferenças que a quarta questão de investigação se constrói. Se a direção é
o sinal fraco que aquele estudo de ablação mostra, e se ali o espaço de comparação permanece
inalterado, então ficam em aberto duas perguntas distintas: se a \emph{magnitude} do movimento,
que não é uma direção e não constitui previsão, é aprendível do texto, e se o ganho
aparece quando se ajusta o próprio codificador em vez de uma rede colocada por cima dele. A
Secção~\ref{sec:met_qi4} descreve como ambas são postas à prova, com um braço treinado sobre a
magnitude e um braço de replicação treinado sobre a direção, e a Secção~\ref{sec:av_qi4}
reporta o que se obteve.

A representação de cada notícia por um vetor converte a identificação de casos comparáveis num
problema de recuperação de informação. A medida de proximidade adotada é a semelhança do cosseno,
cuja mecânica é descrita na Secção~\ref{sec:met_recuperacao}. A operação é exata, uma vez que se
percorre a totalidade da base de casos e se ordena o resultado, pelo que este não depende de
qualquer aproximação. A escalas substancialmente superiores, a procura exaustiva pode ser
substituída por um índice aproximado \autocite{johnson2021faiss} sem alteração da representação, o
que constitui caminho de crescimento e não necessidade presente.

Uma vez que a recuperação devolve uma lista ordenada, a sua avaliação recorre a medidas sensíveis à
ordem \autocite{manning2008ir}. A adotada neste trabalho é a precisão nos $k$ primeiros resultados,
ou seja a fração desses $k$ resultados que são relevantes. A definição de relevância, que é onde uma
avaliação desta natureza se decide, é estabelecida na Secção~\ref{sec:met_recuperacao} juntamente
com as linhas de base contra as quais o resultado é comparado.

\subsection{Recuperação com geração, e por que razão o sistema para antes dela}
\label{sec:ctx_rag}

A arquitetura hoje corrente para responder a uma pergunta a partir de um corpo de documentos é a
geração aumentada por recuperação, na qual um recuperador seleciona passagens e um modelo de
linguagem redige a resposta a partir delas \autocite{lewis2020rag}. A literatura sobre a família é
extensa e encontra-se sistematizada \autocite{huang2026ragsurvey}.

O sistema aqui descrito executa a primeira metade dessa arquitetura e detém-se antes da segunda, e a
razão é a mesma que atravessa o trabalho. A recuperação devolve objetos verificáveis: um título, uma
data, uma semelhança medida e um movimento de preço observado. A geração converte esses objetos numa
afirmação em prosa, cuja fidelidade ao que foi recuperado é ela própria uma propriedade a avaliar, e
que o destinatário deste sistema não tem meios para confrontar. A opção não decorre de limitação
técnica, e a distinção deve ser feita com clareza: acrescentar a geração é possível e barato; o que
não é barato é a avaliação de que ela não desloca o significado da evidência. Nenhuma comparação com
uma arquitetura completa foi realizada neste trabalho, pelo que o que aqui se afirma é uma razão de
desenho e não um resultado. A Secção~\ref{sec:sis_ciclo} descreve a camada de geração ancorada que
chegou a ser construída e avaliada, e a razão pela qual não é exposta.

\subsection{Corpora de notícias financeiras e rótulos de setor}

A avaliação de recuperação sobre notícias financeiras exige um corpo de texto alinhado com preços, e
esse alinhamento é o que distingue um corpus utilizável de uma coleção de títulos. O conjunto
utilizado neste trabalho é o \gls{FNSPID} \autocite{dong2024fnspid}, que reúne notícias e séries de
preços correspondentes e cuja construção está documentada.

A relevância entre duas notícias é aproximada, neste trabalho, pela pertença ao mesmo setor. A opção substitui uma anotação humana inexistente por um critério objetivo, e transporta a
limitação correspondente. A validade de um agrupamento é mensurável, pela separação entre grupos
\autocite{rousseeuw1987silhouettes} ou pela concordância entre agrupamentos alternativos corrigida
para o acaso \autocite{vinh2010ami}. Nenhuma dessas medidas foi aplicada ao mapa setorial adotado. O que o rótulo garante é objetividade e reprodutibilidade, e não que o setor seja
a partição mais informativa do corpus.

\section{Estudos de evento e decomposição de retornos}
\label{sec:ctx_evento}

Esta secção apresenta o instrumento com que se mede o efeito de um acontecimento sobre o preço, bem
como a correção estatística de que depende a separação entre a componente de empresa e a componente
de mercado.

O instrumento adotado é o estudo de evento, cuja formulação remonta a
\textcite{fama1969adjustment}. O comportamento do método sobre retornos diários foi examinado por
\textcite{brown1985daily}, que confrontam formas alternativas de estimar o retorno esperado e
concluem que os procedimentos correntes, assentes na estimação do modelo de mercado por mínimos
quadrados, permanecem bem especificados em janelas curtas em torno do acontecimento. O alcance
desta conclusão é delimitado pelos próprios autores: a variante mais simples, que subtrai apenas a
média histórica da ação, perde potência estatística e pode tornar-se mal especificada quando as
datas dos acontecimentos coincidem. É este resultado que motiva a ordem de grandeza das
janelas de um, três e cinco dias adotadas neste trabalho. A transposição não é direta, e a diferença
está declarada na
Secção~\ref{sec:met_recuperacao}, segundo a qual o valor apresentado ao utilizador é o retorno bruto
e não o retorno anormal que aqueles autores estudam.

A existência de uma relação estatística entre o texto publicado e os movimentos do mercado é
observação antiga \autocite{tetlock2007media}. O verbo adequado é o da evidência de uma relação
estatística, e não o da demonstração de causalidade.

Um ponto metodológico condiciona a segunda pergunta do trabalho, a da separação entre a empresa e o
mercado. Essa separação exige estimar a sensibilidade de cada ação ao mercado, e a estimativa é
ruidosa, sobretudo em empresas voláteis ou com histórico curto. \textcite{blume1971risk} estabelece
que os coeficientes assim estimados tendem a regredir para a média entre períodos sucessivos, do
que decorreu a prática de encolher a estimativa com um peso fixo. Este trabalho adota, em
substituição, o encolhimento ponderado pela imprecisão da própria estimativa \autocite{vasicek1973beta}, cuja
formulação e efeito são apresentados na Secção~\ref{sec:met_decomposicao}.

O enquadramento teórico que justifica a recusa de prever pertence igualmente a esta área. A
hipótese dos mercados eficientes, na sua forma semiforte, sustenta que os preços refletem já a
informação publicamente disponível \autocite{fama1970efficient}, do que decorre que movimentos
futuros não deveriam ser sistematicamente previsíveis a partir dessa informação. Constitui uma das
razões pelas quais o sistema mede o que já ocorreu em vez de projetar o que virá; a segunda, de
natureza prática, é que apenas o passado é verificável por quem lê.

\section{Raciocínio baseado em casos}
\label{sec:ctx_casos}

Esta secção enquadra a terceira técnica do sistema numa linha estabelecida de inteligência
artificial e delimita a parte dessa linha que o trabalho implementa.

O raciocínio baseado em casos \autocite{aamodt1994cbr} aborda um problema novo através da
recuperação de problemas anteriores semelhantes e do aproveitamento do que sobre eles é conhecido,
dispensando o recurso exclusivo a conhecimento geral do domínio. A fonte organiza o método num
ciclo de quatro processos, a saber recuperar, reutilizar, rever e reter, e considera os quatro
necessários à resolução completa de um problema.

A Figura~\ref{fig:ctx_cbr} representa o ciclo e assinala onde o sistema se detém.

\begin{figure}[ht]
\centering
\begin{tikzpicture}[font=\footnotesize,
  cheio/.style={draw=black!75, line width=0.9pt, rounded corners=2pt, fill=black!6,
            minimum width=2.45cm, minimum height=1.05cm, align=center},
  vazio/.style={draw=black!45, densely dashed, line width=0.8pt, rounded corners=2pt,
            fill=white, minimum width=2.45cm, minimum height=1.05cm, align=center,
            text=black!55},
  leg/.style={font=\scriptsize, text=black!65, align=center},
  seta/.style={-{Stealth[length=4pt]}, draw=black!50},
]
  \node[cheio] (n0) at (0.0,0) {\textbf{recuperar}\\[1pt]\scriptsize encontra casos\\ semelhantes};
  \node[cheio] (n1) at (3.15,0) {\textbf{reutilizar}\\[1pt]\scriptsize aproveita o que\\ neles foi medido};
  \node[vazio] (n2) at (6.3,0) {\textbf{rever}\\[1pt]\scriptsize corrige o desfecho\\ para o caso novo};
  \node[vazio] (n3) at (9.45,0) {\textbf{reter}\\[1pt]\scriptsize guarda a solução\\ como caso futuro};
  \draw[seta] (n0) -- (n1);
  \draw[seta] (n1) -- (n2);
  \draw[seta] (n2) -- (n3);
  \draw[seta, densely dashed, draw=black!35] (n3) -- ++(0,-0.95) -- node[leg, below]{} (0,-0.95) -- (n0);
  \node[leg] at (1.58,0.95) {o sistema implementa};
  \node[leg] at (7.87,0.95) {o sistema para aqui};
  \draw[draw=black!30] (4.72,1.15) -- (4.72,-0.62);
  \node[leg, text width=8.2cm] at (4.72,-1.7) {porque adaptar produziria uma expectativa sobre o que vem a seguir, que é a previsão que o trabalho recusa};
\end{tikzpicture}
\caption[O ciclo do raciocínio baseado em casos, e onde o sistema para]{Os quatro processos que \textcite{aamodt1994cbr} consideram necessários à resolução completa de um problema. O sistema executa os dois primeiros a cheio e detém-se antes dos dois seguintes. A interrupção não é uma insuficiência de implementação: \emph{rever} adapta o desfecho de um caso passado ao caso presente, e essa adaptação é uma expectativa sobre o futuro, ou seja exatamente a previsão que a restrição fundadora recusa. O que o sistema apresenta é o desfecho medido de cada caso, tal como ocorreu. Os dois primeiros processos, executados por inteiro e com a decisão entregue à pessoa, correspondem ao que \textcite{kolodner1991aiding} designa por apoio à decisão baseado em casos.}
\label{fig:ctx_cbr}
\end{figure}
A técnica descrita na Secção~\ref{sec:met_recuperacao} implementa os dois primeiros e interrompe-se
deliberadamente antes dos restantes: recupera casos semelhantes e apresenta o que neles foi medido,
sem adaptar esse desfecho ao caso presente. A adaptação produziria uma expectativa sobre o que virá
a acontecer, ou seja precisamente a previsão que o trabalho recusa. A implementação parcial de um
ciclo estabelecido é, neste caso, uma consequência da restrição fundadora e não uma insuficiência.

Esta paragem tem, além disso, nome próprio na literatura, e reconhecê-lo altera-lhe o estatuto.
\textcite{kolodner1991aiding} descreve o \emph{apoio à decisão baseado em casos} como uma
metodologia em que o sistema aumenta a memória da pessoa fornecendo-lhe casos análogos, e em que é
a pessoa, e não a máquina, quem decide, servindo-se desses casos como orientação. O ponto de
partida é uma observação de psicologia: as pessoas raciocinam bem por analogia mas recordam mal os
casos certos, ao passo que os computadores recordam bem. Enunciado assim, o que este sistema faz
deixa de ser a execução de metade de um ciclo e passa a ser a execução completa de uma metodologia
distinta, proposta em 1991 por quem definiu a área. A diferença não é de amor-próprio: uma
implementação parcial explica o que lhe falta, ao passo que uma metodologia declara o que entrega
e a quem cabe decidir.

Convém, em sentido contrário, registar o que o raciocínio baseado em casos tem efetivamente feito
neste domínio, para que a originalidade reivindicada não seja maior do que é.
\textcite{oh2007cbrmonitoring} constroem um indicador diário da condição do mercado financeiro
coreano que classifica cada dia como período estável, instável ou de crise, recuperando casos de
uma base construída sobre a crise de 1997. É raciocínio baseado em casos aplicado a finanças, e é
instrutivo por dois motivos. O primeiro é o objeto: a unidade classificada é o mercado inteiro,
num juízo sobre risco sistémico, e não o movimento de uma empresa na sequência de uma notícia. O
segundo é o destino dos casos recuperados, que servem para produzir automaticamente uma etiqueta;
o artigo não descreve a sua apresentação a quem quer que seja. Mostrar os casos, que é o que aqui
se faz com eles, é precisamente o que ali fica por fazer.

A opção encontra apoio adicional do lado de quem lê. A investigação sobre explicação em ciências
sociais estabelece que as pessoas procuram razões contrastivas e seletivas, e não uma enumeração
exaustiva de tudo o que contribuiu para um resultado \autocite{miller2019explanation}. Um conjunto
reduzido de casos concretos, com data e desfecho, corresponde melhor a essa expectativa do que um
coeficiente. É por essa razão que cada alerta apresenta alguns precedentes e as quantidades exatas
que sustentam a decisão, em vez da totalidade dos valores calculados.

Há um segundo requisito, do mesmo lado, que condiciona não o conteúdo da explicação mas a frequência
com que ela é entregue. Um destinatário exposto a avisos em excesso deixa de os processar, efeito medido sobre
investidores particulares \autocite{liu2023alerts} e discutido na
Secção~\ref{sec:ctx_investidor}. A consequência é que a política de decisão constitui parte do desenho. Um alerta que não
é lido vale o mesmo que um alerta que não foi enviado, e é por isso que o
Capítulo~\ref{cap:sistema} trata os critérios de supressão com o cuidado que dedica aos métodos.

\section{Explicabilidade e confiança em sistemas automáticos}
\label{sec:ctx_xai}

Esta secção situa o requisito de explicabilidade que atravessa o trabalho e apresenta as cautelas
que a literatura associa à apresentação de explicações.

A exigência de que um sistema demonstre como chegou a uma conclusão constitui uma área própria, com
revisões consolidadas \autocite{arrieta2020xai, adadi2018peeking}. A área divide-se em duas vias,
representadas na Figura~\ref{fig:ctx_xai}: os modelos transparentes, interpretáveis em si mesmos, e
os métodos \emph{a posteriori}, que explicam um modelo opaco já treinado e que podem ser
organizados pelo tipo de objeto que explicam \autocite{guidotti2018survey}.

\begin{figure}[ht]
\centering
\begin{tikzpicture}[
  font=\small, >={Stealth[length=2.4mm]},
  r/.style={draw, rounded corners, align=center, text width=3.5cm, minimum height=1.1cm,
            inner sep=4pt, fill=black!3},
  t/.style={draw, rounded corners, align=center, text width=3.5cm, minimum height=1.4cm,
            inner sep=4pt, fill=black!8, very thick},
  o/.style={draw, rounded corners, align=center, text width=3.5cm, minimum height=1.4cm,
            inner sep=4pt, fill=black!3},
]
  \node[r] (raiz) {\textbf{Explicar uma decisão}};
  \node[t, below left=1.1cm and 0.35cm of raiz] (tr) {\textbf{Modelo transparente}\\
    \footnotesize a explicação é o próprio cálculo};
  \node[o, below right=1.1cm and 0.35cm of raiz] (ph) {\textbf{Método \emph{a posteriori}}\\
    \footnotesize explica um modelo opaco};

  \node[below=0.4cm of tr, align=center, font=\footnotesize, text width=3.6cm]
    (trx) {regra, regressão linear,\\ contagem sobre a série};
  \node[below=0.4cm of ph, align=center, font=\footnotesize, text width=3.6cm]
    (phx) {\gls{LIME}, \gls{SHAP}};

  \draw[->, thick, black!45] (raiz) -- (tr);
  \draw[->, thick, black!45] (raiz) -- (ph);
  \draw[-, black!35] (tr) -- (trx);
  \draw[-, black!35] (ph) -- (phx);
\end{tikzpicture}
\caption[As duas vias para a explicabilidade]{As duas vias distinguidas pela literatura
\autocite{arrieta2020xai, guidotti2018survey}: construir modelos transparentes por desenho, ou
explicar um modelo opaco depois de treinado. Este trabalho segue a primeira.}
\label{fig:ctx_xai}
\end{figure}

Os dois métodos \emph{a posteriori} de utilização mais generalizada operam de formas distintas. O
\gls{LIME} ajusta um modelo simples na vizinhança de uma previsão individual, aproximando
localmente o comportamento do modelo opaco \autocite{ribeiro2016lime}. O \gls{SHAP} atribui uma
previsão às variáveis que a produziram recorrendo a valores de Shapley \autocite{lundberg2017shap}.
Ambos são genéricos e ambos são úteis.

Ambos partilham, contudo, uma propriedade determinante neste contexto. Perante decisões de
consequência séria, \textcite{rudin2019stop} defende o abandono dos modelos opacos em favor de
modelos interpretáveis por construção, com o argumento de que toda a explicação \emph{a posteriori}
constitui ela própria uma aproximação, e uma aproximação pode conduzir a leituras erradas. O
utilizador recebe nesse caso dois objetos cuja relação não está em condições de avaliar, que são
uma decisão e uma explicação que pode não lhe corresponder.

Este trabalho segue a via conservadora. Em lugar de explicar \emph{a posteriori} um modelo opaco,
recorre a componentes em que a explicação e o cálculo coincidem. Um afastamento face a uma norma
recente, uma repartição aditiva cuja soma reproduz o movimento observado e um conjunto de casos
anteriores com os desfechos medidos constituem objetos apresentáveis por inteiro. A explicação
não constitui um segundo objeto produzido sobre a decisão: é a decisão, enunciada por extenso.

Duas cautelas da literatura fecham esta secção, e ambas condicionam a organização do
Capítulo~\ref{cap:avaliacao}.

A primeira é que interpretabilidade não designa uma propriedade única \autocite{lipton2018mythos},
pelo que um sistema deve declarar o tipo de interpretabilidade que oferece em vez de reclamar o
termo sem qualificação. Este trabalho oferece transparência de mecanismo, ou seja o utilizador pode
seguir o cálculo que conduziu à decisão, e não oferece uma explicação causal do acontecimento no
mundo.

A segunda respeita ao efeito da própria explicação sobre quem a recebe. A confiança em sistemas
automáticos falha em dois sentidos simétricos, com o utilizador a ignorar um aviso correto ou a
aceitar um aviso incorreto, pelo que o objetivo não é maximizar a confiança mas torná-la ajustada à
capacidade real do sistema \autocite{lee2004trust}. \textcite{bansal2021whole} acrescentam que a
presença de uma explicação aumenta a probabilidade de aceitação da resposta do sistema quer esta
esteja certa quer esteja errada, ou seja que o ganho observado não decorre de a pessoa passar a
distinguir melhor os dois casos. As consequências de desenho que daqui decorrem estão descritas na
Secção~\ref{sec:met_etica}.

Estas cautelas são gerais. Sobre a opção concreta desta dissertação, que é explicar por casos,
existe evidência empírica recolhida com pessoas que aponta em sentido contrário, e remetê-la para
uma nota de rodapé seria a forma mais barata de a neutralizar.

\textcite{waa2021xai} compararam explicações por regras, explicações por exemplos e ausência de
explicação, em duas experiências com quarenta e cinco participantes cada, num apoio à decisão para
dosagem de insulina. As explicações por regras melhoraram significativamente a identificação do
fator decisivo; as explicações por exemplos não se distinguiram da ausência de explicação, nem na
identificação do fator ($p = 0{,}796$) nem na compreensão auto-reportada ($p = 0{,}283$). A
explicação que os autores dão para o resultado é a parte que mais interessa a este trabalho:
ambos os estilos, escrevem, fornecem apenas os pormenores relevantes para uma decisão isolada, e
não o raciocínio ou a causalidade que lhe estão por baixo.

No próprio domínio financeiro, e com não-especialistas, existem dois estudos recentes.
\textcite{cau2023logicstyle} comparam três estilos de explicação numa tarefa de negociação de
ações e concluem que os estilos abdutivo e dedutivo melhoram o desempenho quando o modelo está
confiante, por comparação com o estilo indutivo, que é o estilo baseado em exemplos e portanto o
mais próximo daquilo que este sistema entrega. \textcite{bertrand2023featurebased} testam um
consultor automático de seguros de vida com explicações por atributos e concluem que estas não
melhoram nem a compreensão nem a adequação da confiança, por comparação com não dar explicação
nenhuma, e que o formato mais conversacional aumenta a confiança do utilizador por vezes em
prejuízo dele.

Daqui decorrem três consequências para esta dissertação. A primeira é que a opção por explicação
baseada em casos, justificada na Secção~\ref{sec:ctx_casos} pela literatura de explicação em
ciências sociais \autocite{miller2019explanation}, não tem apoio nas medições com pessoas que
existem, e num caso tem contra si evidência do próprio domínio. A segunda é que o diagnóstico de
\textcite{waa2021xai} nomeia exatamente aquilo que este sistema acrescenta aos casos: uma
repartição aditiva cuja soma reproduz o movimento observado, que é o cálculo subjacente cuja
ausência os autores responsabilizam pelo resultado. O sistema não entrega apenas casos; entrega
casos e a conta que os acompanha. Se isso é suficiente é uma pergunta empírica em aberto, e o
trabalho não a responde. A terceira, e a mais consequente, é que estes resultados tornam o estudo
com utilizadores mais necessário e não menos: existe literatura publicada a apontar contra a
hipótese fundadora, e é nesses termos que a Secção~\ref{sec:con_limitacoes} enuncia a sua
ausência.

O segundo estudo apoia, em contrapartida, uma decisão já tomada. Se o formato conversacional é o
que produz confiança excessiva, então não expor a camada conversacional descrita na
Secção~\ref{sec:sis_inteligencia} deixa de ser apenas prudência de engenharia e passa a ter apoio
empírico no domínio.

Acresce que afirmações de interpretabilidade exigem avaliação rigorosa e não se estabelecem por
declaração \autocite{doshivelez2018considerations}. É por essa razão que o Capítulo~\ref{cap:avaliacao}
mede o que é mensurável por procedimento automático e identifica explicitamente a parte que
permanece por avaliar por ausência de um estudo controlado com utilizadores.

\section{Aprendizagem automática em ambiente de produção}
\label{sec:ctx_producao}

Esta secção reúne a literatura relativa ao componente aprendido do sistema, organizada segundo as
três questões que a sua implantação levanta.

A primeira questão é a de determinar contra que referência superior comparar um modelo simples.
Conjuntos de árvores treinados por reforço do gradiente \autocite{friedman2001gbm} captam interações
e efeitos não lineares que uma regressão logística não capta, e desempenham por isso essa função no
Capítulo~\ref{cap:avaliacao}. Não se afirma que constituam o melhor método existente, uma vez que
tal seria uma afirmação sobre um campo inteiro.

A segunda questão respeita à honestidade do valor apresentado. Uma pontuação só é interpretável
como probabilidade se estiver calibrada, e as pontuações em bruto de muitos classificadores não o
estão; a calibração \emph{a posteriori} sobre um bloco de validação reservado corrige essa
distorção \autocite{niculescu2005calibration}. A mesma fonte não favorece, contudo, a opção tomada
neste trabalho: identifica a regressão logística como já bem calibrada à partida e conclui que,
acima de cerca de mil casos de calibração, a regressão isotónica iguala ou supera a sigmoide. O
modelo implantado é uma regressão logística com um bloco de validação muito acima desse limiar,
pelo que ambas as afirmações se lhe aplicam. A opção justifica-se por medição e não por autoridade:
a comparação das duas variantes sobre a mesma validação está reportada na
Tabela~\ref{tab:av_alternativas}, na página~\pageref{tab:av_alternativas}.

Uma via alternativa foi considerada e não adotada. A predição conformal devolve um conjunto de
rótulos com uma garantia livre de distribuição, válida em amostras finitas
\autocite{vovk2005algorithmic, angelopoulos2023conformal}. Duas precisões determinam que não é aplicável nas condições deste trabalho. A cobertura da
variante corrente é marginal: vale em média sobre os casos, e não caso a caso. E a garantia exige
permutabilidade entre o bloco de calibração e o de teste, hipótese que a divisão cronológica com
embargo da Secção~\ref{sec:met_triagem} recusa por construção. A Secção~\ref{sec:con_futuro} retoma o método
como caminho a explorar.

A terceira questão é a da permanência do desempenho ao longo do tempo. Um modelo implantado enfrenta
alteração da distribuição dos dados que observa, e a literatura de deriva de conceito estabelece a
taxonomia dessa mudança e as estratégias de deteção e adaptação \autocite{gama2014survey}. Este
trabalho mede a deriva sobre os seus próprios dados em vez de a assumir ausente.

Existe, por fim, um custo que não é captado por nenhuma métrica de qualidade.
\textcite{sculley2015debt} estabelecem que os componentes de aprendizagem automática acumulam
manutenção oculta, em particular através do emaranhamento entre peças e de dependências de dados
instáveis e não declaradas. O Capítulo~\ref{cap:conclusoes} regressa a este ponto com um exemplo
medido no presente sistema, em que parte do ciclo de aprendizagem se interrompeu sem que qualquer
erro fosse registado.

\section{Síntese e posicionamento do trabalho}
\label{sec:ctx_sintese}

Esta secção localiza o trabalho relativamente à revisão apresentada e enuncia as decisões de âmbito
que dela decorrem.

A revisão permite delimitar com precisão o espaço em que o trabalho se inscreve, e esse espaço não
constitui um vazio técnico. Cada componente utilizado existe, está publicado e está avaliado, como
a Tabela~\ref{tab:ctx_posicionamento} torna explícito.

\begin{table}[!htbp]
\caption[Posicionamento do trabalho face à literatura]{Cada componente do sistema existe na
literatura de forma isolada. O que não existe é a sua combinação sob a restrição de custo nulo, com
a evidência entregue ao utilizador no momento em que a afirmação lhe é apresentada.}
\label{tab:ctx_posicionamento}
\centering
\small
\begin{tabular}{L{0.24\textwidth} L{0.30\textwidth} L{0.34\textwidth}}
\toprule
\tabhead{Pergunta} & \tabhead{O que a literatura oferece} & \tabhead{O que falta ao público-alvo} \\
\midrule
É invulgar?
  & Deteção de anomalias, estatística e aprendida \autocite{chandola2009anomaly, liu2008isolation}
  & Uma medida relativa à norma da própria empresa, comunicada em linguagem acessível \\
\addlinespace
É a empresa ou o mercado?
  & Estudo de evento e estimação de sensibilidade \autocite{brown1985daily, vasicek1973beta}
  & Disponibilidade fora de um terminal profissional \\
\addlinespace
Já aconteceu antes?
  & Representação de frases e recuperação \autocite{reimers2019sbert, manning2008ir}
  & Uma base de casos com desfechos medidos, e não recordados \\
\addlinespace
Por que razão acreditar?
  & Explicabilidade e literatura de confiança \autocite{rudin2019stop, lee2004trust}; medições
    com pessoas que não confirmam o ganho \autocite{waa2021xai, cau2023logicstyle,
    bertrand2023featurebased}
  & A evidência entregue junto da afirmação e verificável no momento da leitura; por
    demonstrar com utilizadores \\
\bottomrule
\end{tabular}
\end{table}

A lacuna não é, portanto, de método, mas de integração sob restrição. Os produtos que respondem às três perguntas são terminais profissionais, indisponíveis em regime
gratuito ou acessível ao investidor particular. O seu preço não é publicado de forma citável, e é
a indisponibilidade, e não um valor concreto, que sustenta o argumento. Os produtos
gratuitos que declaram responder à terceira entregam uma afirmação sem a evidência que a sustenta.
Na literatura académica há sistemas integrados e propostas de explicação de movimentos passados,
como a Secção~\ref{sec:ctx_ferramentas} documenta. A distinção face ao MoFiN DSS inclui a ausência
de previsão; face ao DeepTrust e ao Stock Pattern Assistant, essa opção é partilhada. A
contribuição deve, por isso, ser apreciada pela articulação entre raridade, decomposição e
precedentes, pelas restrições de operação e pela avaliação da evidência entregue. Os trabalhos
revistos delimitam esta contribuição; não demonstram uma ausência universal de sistemas comparáveis.

Esta observação determina a distinção que atravessa o trabalho e que o Capítulo~\ref{cap:sistema}
retoma ao justificar a ausência de um modelo de linguagem na composição das mensagens. Uma frase
gerada sem ligação aos factos apresenta-se ao leitor num registo fluente e não lhe fornece meios para a
confrontar com aquilo que a sustenta. Uma frase composta a partir de objetos calculados é uma
afirmação acompanhada da evidência, na qual cada valor apresentado remete para a medição que o
produziu e o leitor pode confirmar cada elemento no momento em que o lê. É esta a propriedade que o sistema procura preservar, e é ela que fixa a natureza da contribuição.
Não é a formulação de um método novo. É a seleção, integração e avaliação crítica de métodos
existentes num sistema funcional, sob restrições declaradas, e com os resultados que não favorecem
as opções tomadas incluídos na avaliação.

Quatro decisões de âmbito decorrem deste posicionamento, e cada uma exclui deliberadamente uma
classe de funcionalidade.

\begin{description}
  \item[Responder às três perguntas, e apenas a essas.] Não são incorporadas funcionalidades que
        não sirvam uma delas, o que mantém o sistema com dimensão compatível com a sua explicação
        integral.
  \item[Não produzir previsões de preço.] Nem direção, nem alvo, nem recomendação. Além da razão de
        verificabilidade enunciada no Capítulo~\ref{cap:introducao}, aplica-se o enquadramento
        teórico da Secção~\ref{sec:ctx_evento} \autocite{fama1970efficient}.
  \item[Não gerir carteiras nem prestar aconselhamento.] O sistema não conhece as posições do
        utilizador e não lhe indica que ação tomar, o que evita a recolha de dados pessoais e
        mantém o trabalho fora da fronteira do aconselhamento financeiro regulado.
  \item[Utilizar exclusivamente fontes gratuitas.] A restrição torna o sistema acessível ao público
        a que se destina e converte as limitações resultantes em quantidades medidas e reportadas.
\end{description}

A última decisão tem um custo real que importa enunciar sem atenuação. As notícias chegam com
atraso e a cobertura dos acontecimentos é incompleta. O Capítulo~\ref{cap:sistema} quantifica
ambas as limitações e o Capítulo~\ref{cap:conclusoes} avalia as suas consequências.
